\section{The pressure-complete setting and mixed jet}
\label{f:sec:formal-setting-mixed-jet}

\begin{proposition}[Pressure normalization]
\label{f:prop:formal-pressure-normalization}
Let \(u,p,f\) be smooth on \(I\times\R^3\), let \(\nabla\cdot u=0\), and suppose
\[
 u(t)\in H^1_\sigma(\R^3),
 \qquad
 p(t)\in H^0(\R^3),
 \qquad f(t)\in L^1(\R^3)\cap H^1(\R^3)
 \qquad(t\in I),
\]
and
\begin{equation}
 u_t+(u\cdot\nabla)u+\nabla p=\nu\Delta u+f.
 \label{f:eq:formal-ns-system}
\end{equation}
Then
\begin{align}
 -\Delta p
 &=\partial_i\partial_j(u_i u_j)-\nabla\cdot f,
 \label{f:eq:formal-pressure-poisson}\\
 p
 &=R_iR_j(u_i u_j)-(-\Delta)^{-1}\nabla\cdot f,
 \label{f:eq:formal-pressure-riesz}\\
 (u\cdot\nabla)u+\nabla p
 &=\Pp\nabla\cdot(u\otimes u)+(I-\Pp)f.
 \label{f:eq:formal-pressure-projected}
\end{align}
The condition \(p(t)\in H^0(\R^3)\) fixes the pressure uniquely.
\end{proposition}

\begin{proof}
Incompressibility gives
\[
 (u\cdot\nabla)u_i=\partial_j(u_j u_i).
\]
Taking the divergence of \eqref{f:eq:formal-ns-system} cancels
\(\nabla\cdot u_t\) and \(\nabla\cdot\Delta u\), which yields
\eqref{f:eq:formal-pressure-poisson}.  Lemma~\ref{f:lem:riesz-leray}, applied
with \(m=0\), gives the force potential in \(H^0\) and
\[
 q:=R_iR_j(u_i u_j)-(-\Delta)^{-1}\nabla\cdot f\in H^0.
\]
It satisfies this Poisson identity and
\[
 (u\cdot\nabla)u+\nabla q
 =\Pp\nabla\cdot(u\otimes u)+(I-\Pp)f.
\]
If \(p_1,p_2\in H^0\) satisfy \eqref{f:eq:formal-pressure-poisson}, then
\(h=p_1-p_2\) satisfies
\(|\xi|^2\widehat h=0\).  Since \(\widehat h\in L^2\) and is supported at
\(\{0\}\), it vanishes.  Hence \(p=q\), which proves
\eqref{f:eq:formal-pressure-riesz} and \eqref{f:eq:formal-pressure-projected}.
\end{proof}

\begin{samepage}
\begin{definition}[Maximal pressure-complete classical history]
\label{f:def:formal-maximal-history}
Fix \(\nu>0\), \(u_0\in\HsSigma(\R^3)\), and
\(f\in C^\infty([0,\infty);\mathcal S(\R^3;\R^3))\).  A
pressure-complete classical history generated by \((u_0,\nu,f)\) is a triple
\[
 (T_*,u,p),
 \qquad
 0<T_*\le\infty,
\]
such that, for every \(0<T<T_*\) and \(n,m\in\Nzero\),
\begin{align*}
 \partial_t^nu&\in C([0,T];H^m_\sigma(\R^3)),\\
 \partial_t^np&\in C([0,T];H^m(\R^3)),
\end{align*}
and
\begin{align}
 u_t+(u\cdot\nabla)u+\nabla p&=\nu\Delta u+f,
 \label{f:eq:formal-history-momentum}\\
 \nabla\cdot u&=0,
 \label{f:eq:formal-history-divergence}\\
 u(0)&=u_0,
 \label{f:eq:formal-history-datum}\\
 p(t)&=R_iR_j(u_i(t)u_j(t))
 -(-\Delta)^{-1}\nabla\cdot f(t).
 \label{f:eq:formal-history-pressure}
\end{align}
The history is maximal if no pressure-complete classical history generated by
the same \((u_0,\nu,f)\) continues \((u,p)\) past \(T_*\).
\end{definition}
\end{samepage}

\begin{definition}[Temporal and mixed velocity--pressure jets]
\label{f:def:formal-mixed-jet}
Let \((T_*,u,p)\) be a maximal pressure-complete classical history.  For
\(n\in\Nzero\) and \(\alpha\in\Nzero^3\), define
\[
 V_n:=\partial_t^nu,
 \qquad
 Q_n:=\partial_t^np,
 \qquad
 F_n:=\partial_t^nf,
 \qquad
 J_{n,\alpha}:=\partial_t^n\partial_x^\alpha u,
 \qquad
 \Pi_{n,\alpha}:=\partial_t^n\partial_x^\alpha p.
\]
Set also \(F_{n,\alpha}:=\partial_t^n\partial_x^\alpha f\).
For \(\beta\le\alpha\), define
\[
 \binom{\alpha}{\beta}
 :=\prod_{\ell=1}^3\binom{\alpha_\ell}{\beta_\ell},
 \qquad
 |\alpha|:=\alpha_1+\alpha_2+\alpha_3.
\]
\end{definition}

\begin{proposition}[Pressure-complete binomial recurrences]
\label{f:prop:formal-mixed-recurrences}
For every \(n\in\Nzero\) and \(\alpha\in\Nzero^3\),
\begin{align}
 Q_n
 &=R_iR_j
 \sum_{a=0}^{n}\binom na V_{a,i}V_{n-a,j}
 -(-\Delta)^{-1}\nabla\cdot F_n,
 \label{f:eq:formal-temporal-pressure-row}\\
 V_{n+1}
 &=
 \nu\Delta V_n
 -\sum_{a=0}^{n}\binom na(V_a\cdot\nabla)V_{n-a}
 -\nabla Q_n+F_n,
 \label{f:eq:formal-temporal-velocity-row}\\
 \Pi_{n,\alpha}
 &=R_iR_j
 \sum_{a=0}^{n}\sum_{\beta\le\alpha}
 \binom na\binom{\alpha}{\beta}
 J_{a,\beta,i}J_{n-a,\alpha-\beta,j}
 -\partial_x^\alpha(-\Delta)^{-1}\nabla\cdot F_n,
 \label{f:eq:formal-mixed-pressure-row}\\
 J_{n+1,\alpha}
 &+
 \sum_{a=0}^{n}\sum_{\beta\le\alpha}
 \binom na\binom{\alpha}{\beta}
 (J_{a,\beta}\cdot\nabla)
 J_{n-a,\alpha-\beta}
 +\nabla\Pi_{n,\alpha}
 =\nu\Delta J_{n,\alpha}+F_{n,\alpha},
 \label{f:eq:formal-mixed-velocity-row}\\
 \nabla\cdot J_{n,\alpha}&=0.
 \label{f:eq:formal-mixed-divergence-row}
\end{align}
\end{proposition}

\begin{proof}
Apply \(\partial_t^n\) to \eqref{f:eq:formal-history-pressure} and use
Lemma~\ref{f:lem:mixed-leibniz} with \(\alpha=0\) to obtain
\eqref{f:eq:formal-temporal-pressure-row}.  Apply \(\partial_t^n\) to
\eqref{f:eq:formal-history-momentum} to obtain
\eqref{f:eq:formal-temporal-velocity-row}.  Apply
\(\partial_x^\alpha\) to both identities, use
Lemma~\ref{f:lem:mixed-leibniz}, and commute the derivatives with \(R_iR_j\)
to obtain \eqref{f:eq:formal-mixed-pressure-row} and
\eqref{f:eq:formal-mixed-velocity-row}.  Differentiation of
\eqref{f:eq:formal-history-divergence} gives
\eqref{f:eq:formal-mixed-divergence-row}.
\end{proof}

\paragraph{The forced datum and its derivative rows.}
The initial velocity, viscosity, and complete prescribed force specify all
inputs to these recurrences.  Pressure is determined by the velocity products
and the gradient part of that same force.  Every temporal and spatial force
derivative remains in its corresponding row.  In the kinetic-energy
calculation of Proposition~\ref{f:prop:kinetic-energy}, transport and pressure
have zero global pairing with the velocity while the external work remains.
At higher Sobolev orders, the full transport--pressure contribution and force
work are retained in Proposition~\ref{f:prop:completed-height-recurrence}.
These are the governing rows and energy identities used below.

\begin{definition}[Canonical pressure constant]
\label{f:def:formal-canonical-bilinear-constants}
For vector fields \(v,w\), set
\[
 \mathcal Q(v,w):=\sum_{i,j=1}^3R_iR_j(v_i w_j).
\]
For \(m\in\Nzero\), define
\[
 C_m^{\rm p}:=
 \sup_{\substack{
 v,w\in H^{m+1}(\R^3;\R^3)\setminus\{0\}}}
 \frac{
 \norm{\mathcal Q(v,w)}_{H^m}
 +\norm{\nabla\mathcal Q(v,w)}_{H^{m-1}}}
 {\norm v_{H^{m+1}}\norm w_{H^{m+1}}}.
\]
\end{definition}

\begin{lemma}
\label{f:lem:formal-canonical-bilinear-constants}
For every \(m\in\Nzero\),
\[
 C_m^{\rm p}\le2C_m^{\rm alg}.
\]
\end{lemma}
\begin{proof}
Regard \((v_i w_j)_{i,j}\) as one tensor.  The contraction
\(F\mapsto R_iR_jF_{ij}\) has Fourier-multiplier norm one from the tensor
\(H^m\) norm to scalar \(H^m\).  Lemma~\ref{f:lem:sobolev-product} and
\(\norm{\nabla h}_{H^{m-1}}\le\norm h_{H^m}\) give
\[
 \norm{\mathcal Q(v,w)}_{H^m}
 +\norm{\nabla\mathcal Q(v,w)}_{H^{m-1}}
 \le2C_m^{\rm alg}\norm v_{H^{m+1}}\norm w_{H^{m+1}}.
\]
\end{proof}

\begin{lemma}[Initial pressure-complete jet]
\label{f:lem:formal-initial-jet}
Set
\[
 V_n^0:=V_n(0),
 \qquad
 Q_n^0:=Q_n(0),
 \qquad
 J_{n,\alpha}^0:=J_{n,\alpha}(0),
 \qquad
 \Pi_{n,\alpha}^0:=\Pi_{n,\alpha}(0).
\]
Write \(F_n^0:=F_n(0)\).
Then
\begin{align}
 V_0^0&=u_0,
 \label{f:eq:formal-initial-base}\\
 Q_n^0
 &=R_iR_j
 \sum_{a=0}^{n}\binom naV_{a,i}^0V_{n-a,j}^0
 -(-\Delta)^{-1}\nabla\cdot F_n^0,
 \label{f:eq:formal-initial-pressure}\\
 V_{n+1}^0
 &=
 \nu\Delta V_n^0
 -\sum_{a=0}^{n}\binom na
 (V_a^0\cdot\nabla)V_{n-a}^0
 -\nabla Q_n^0+F_n^0.
 \label{f:eq:formal-initial-velocity}
\end{align}
For every \(n,m\in\Nzero\) and \(\alpha\in\Nzero^3\),
\[
 V_n^0,Q_n^0,J_{n,\alpha}^0,\Pi_{n,\alpha}^0\in H^m(\R^3).
\]
\par
\end{lemma}

\begin{proof}
Equations \eqref{f:eq:formal-initial-base}--\eqref{f:eq:formal-initial-velocity}
are \eqref{f:eq:formal-temporal-pressure-row} and
\eqref{f:eq:formal-temporal-velocity-row} at \(t=0\).  Assume
\(V_0^0,\ldots,V_n^0\in\bigcap_{m\ge0}H^m\).  Lemmas
\ref{f:lem:sobolev-product} and \ref{f:lem:riesz-leray} give
\[
 Q_n^0\in\bigcap_{m\ge0}H^m,
 \qquad
 V_{n+1}^0\in\bigcap_{m\ge0}H^m.
\]
Induction starts from \(u_0\in\HsSigma\), while every \(F_n^0\) is
Schwartz by the prescribed force class.  Spatial differentiation and
\eqref{f:eq:formal-mixed-pressure-row} give the claims for
\(J_{n,\alpha}^0\) and \(\Pi_{n,\alpha}^0\).
\end{proof}

\begin{theorem}[Finite-row pressure estimates]
\label{f:thm:formal-finite-row-pressure}
Let \(I=(0,T)\subset(0,T_*)\), let \(n,M\in\Nzero\), and suppose
\[
 V_a\in L^2(I;H^{M+1}(\R^3))
 \qquad(0\le a\le n).
\]
For \(\alpha\in\Nzero^3\), define the prescribed pressure-source budget
\begin{equation}
 \mathfrak F^{\rm p}_{n,\alpha,M}(I)
 :=\norm{\partial_x^\alpha(-\Delta)^{-1}\nabla\cdot F_n}_{L^1(I;H^M)}
 +\norm{\nabla\partial_x^\alpha(-\Delta)^{-1}\nabla\cdot F_n}_{L^1(I;H^{M-1})}.
 \label{f:eq:force-pressure-budget}
\end{equation}
It is finite by \eqref{f:eq:force-pressure-potential-bound} and the prescribed
force class.
Then
\begin{align}
 &\norm{Q_n}_{L^1(I;H^M)}
 +\norm{\nabla Q_n}_{L^1(I;H^{M-1})}
 \notag\\
 &\qquad\le
 C_M^{\rm p}\sum_{a=0}^{n}\binom na
 \norm{V_a}_{L^2(I;H^{M+1})}
 \norm{V_{n-a}}_{L^2(I;H^{M+1})}
 +\mathfrak F^{\rm p}_{n,0,M}(I).
 \label{f:eq:formal-finite-pressure-L1}
\end{align}
\begin{samepage}
If
\[
 \mathcal V_{n,M}(I)^2
 :=\sum_{a=0}^{n}
 \norm{V_a}_{L^2(I;H^{M+1})}^2,
\]
then
\begin{equation}
 \norm{Q_n}_{L^1(I;H^M)}
 +\norm{\nabla Q_n}_{L^1(I;H^{M-1})}
 \le C_M^{\rm p}2^n\mathcal V_{n,M}(I)^2
 +\mathfrak F^{\rm p}_{n,0,M}(I).
 \label{f:eq:formal-finite-pressure-compressed}
\end{equation}
\end{samepage}
For \(\alpha\in\Nzero^3\), assume
\[
 J_{a,\beta}\in L^2(I;H^{M+1}(\R^3))
 \qquad
 (0\le a\le n,\ \beta\le\alpha).
\]
Then
\begin{align}
 &\norm{\Pi_{n,\alpha}}_{L^1(I;H^M)}
 +\norm{\nabla\Pi_{n,\alpha}}_{L^1(I;H^{M-1})}
 \notag\\
 &\quad\le
 C_M^{\rm p}
 \sum_{a=0}^{n}\sum_{\beta\le\alpha}
 \binom na\binom{\alpha}{\beta}
 \norm{J_{a,\beta}}_{L^2(I;H^{M+1})}
 \norm{J_{n-a,\alpha-\beta}}_{L^2(I;H^{M+1})}
 +\mathfrak F^{\rm p}_{n,\alpha,M}(I).
 \label{f:eq:formal-finite-mixed-pressure-L1}
\end{align}
If
\[
 \mathcal J_{n,\alpha,M}(I)^2
 :=
 \sum_{a=0}^{n}\sum_{\beta\le\alpha}
 \norm{J_{a,\beta}}_{L^2(I;H^{M+1})}^2,
\]
then
\begin{equation}
 \norm{\Pi_{n,\alpha}}_{L^1(I;H^M)}
 +\norm{\nabla\Pi_{n,\alpha}}_{L^1(I;H^{M-1})}
 \le
 C_M^{\rm p}2^{n+|\alpha|}
 \mathcal J_{n,\alpha,M}(I)^2
 +\mathfrak F^{\rm p}_{n,\alpha,M}(I).
 \label{f:eq:formal-finite-mixed-pressure-compressed}
\end{equation}
\end{theorem}

\begin{proof}
Lemmas~\ref{f:lem:sobolev-product} and \ref{f:lem:riesz-leray} give
\begin{equation}
 \norm{R_iR_j(v_i w_j)}_{H^M}
 +\norm{\nabla R_iR_j(v_i w_j)}_{H^{M-1}}
 \le C_M^{\rm p}\norm v_{H^{M+1}}\norm w_{H^{M+1}}.
 \label{f:eq:formal-pressure-bilinear}
\end{equation}
Equations \eqref{f:eq:formal-temporal-pressure-row} and
\eqref{f:eq:formal-pressure-bilinear} give, for almost every \(t\in I\),
\begin{align*}
 \norm{Q_n(t)}_{H^M}
 +\norm{\nabla Q_n(t)}_{H^{M-1}}
 &\le
 C_M^{\rm p}\sum_{a=0}^{n}\binom na
 \norm{V_a(t)}_{H^{M+1}}
 \norm{V_{n-a}(t)}_{H^{M+1}}\\
 &\quad+\norm{(-\Delta)^{-1}\nabla\cdot F_n(t)}_{H^M}\\
 &\quad+\norm{\nabla(-\Delta)^{-1}\nabla\cdot F_n(t)}_{H^{M-1}}.
\end{align*}
Integration and Cauchy--Schwarz prove
\eqref{f:eq:formal-finite-pressure-L1}, with the two force-potential terms
collected exactly in \(\mathfrak F^{\rm p}_{n,0,M}(I)\).  For nonnegative \(x_a\),
\[
 \sum_{a=0}^{n}\binom nax_ax_{n-a}
 \le
 \sum_{a=0}^{n}\binom nax_a^2
 \le
 2^n\sum_{a=0}^{n}x_a^2.
\]
This proves \eqref{f:eq:formal-finite-pressure-compressed}.

Apply \eqref{f:eq:formal-pressure-bilinear} to each summand in
\eqref{f:eq:formal-mixed-pressure-row}.  Integration and Cauchy--Schwarz give
\eqref{f:eq:formal-finite-mixed-pressure-L1}; the differentiated prescribed
pressure source is \(\mathfrak F^{\rm p}_{n,\alpha,M}(I)\).  The involution
\[
 (a,\beta)\longmapsto(n-a,\alpha-\beta)
\]
preserves the weights
\(\binom na\binom{\alpha}{\beta}\), and
\[
 \sum_{a=0}^{n}\sum_{\beta\le\alpha}
 \binom na\binom{\alpha}{\beta}
 =2^{n+|\alpha|}.
\]
In particular, every individual weight is at most \(2^{n+|\alpha|}\).
The inequality \(2xy\le x^2+y^2\) proves
\eqref{f:eq:formal-finite-mixed-pressure-compressed}.
\end{proof}


\subsection{One compatible forced rectangle family and its two descriptions}
\label{f:sec:independent-intersection}

Start with \(u_0\in\HsSigma(\R^3)\), \(\nu>0\), and the prescribed
\(f\in C^\infty([0,\infty);\mathcal S(\R^3;\R^3))\).
Theorem~\ref{f:thm:local-restart} constructs a pressure-normalized classical
history on a positive interval. Uniqueness identifies the local solutions
on their overlaps; their compatible union gives the maximal history of
Definition~\ref{f:def:formal-maximal-history}.

The datum fixes the initial velocity row. The pressure formula in
Proposition~\ref{f:prop:formal-mixed-recurrences} fixes its normalized pressure,
and the velocity recurrence fixes the next temporal row. Repetition and
spatial differentiation produce the complete initial mixed jet. Along the
history these are the actual derivatives of \(u,p,f\), with the complete
transport sum, pressure, viscosity, and force row retained. These finite jet
realizations give every finite mixed Sobolev row on every strict horizon.

At a proposed finite maximal time, fix finite temporal and spatial orders and
close that index set under the adjacent temporal row, the viscous derivatives,
every Leibniz allocation, the normalized pressure coordinate, and the
prescribed force row.  The datum construction produces that finite
pressure-complete record on \(I=(0,T_*)\), and coordinate deletion recovers
every smaller record.  Ranging through the mixed indices exhausts every
finite temporal and spatial order and gives the direct all-orders mixed
Sobolev membership.  The next section independently derives the
completed-height and viscous identities and evaluates them on these same
coordinates.

\begin{theorem}[Datum-generated strict-horizon mixed-jet intersection]
\label{f:thm:strict-datum-intersection}
Let \(\nu>0\), \(u_0\in\HsSigma(\R^3)\), and
\(f\in C^\infty([0,\infty);\mathcal S(\R^3;\R^3))\). Let
\((T_*,u,p)\) be the maximal pressure-complete classical history generated
by \((u_0,\nu,f)\). For every \(0<T<T_*\),
\begin{equation}
 u\in\bigcap_{r,m\in\Nzero}H^r((0,T);H^m_\sigma(\R^3)).
 \label{f:eq:strict-independent-mixed-intersection}
\end{equation}
All velocity, pressure, and force coordinates are the actual mixed
derivatives of this one history, obey
Proposition~\ref{f:prop:formal-mixed-recurrences}, and agree under restriction.
\end{theorem}

\begin{proof}
Write \(A_n=V_n(0)\) and \(P_n=Q_n(0)\). The complete initial recurrence is
\[
 \begin{aligned}
 A_0&=u_0,\\
 P_n&=R_iR_j\sum_{a=0}^n\binom na(A_a)_i(A_{n-a})_j
      -(-\Delta)^{-1}\nabla\cdot F_n(0),\\
 A_{n+1}&=\nu\Delta A_n
 -\sum_{a=0}^n\binom na(A_a\cdot\nabla)A_{n-a}
 -\nabla P_n+F_n(0).
 \end{aligned}
\]
Lemma~\ref{f:lem:formal-initial-jet} shows inductively that every selected
velocity and pressure row lies in every finite Sobolev space.

Fix \(T<T_*\) and finite \(r,m\). Apply the forced local construction
simultaneously to the finite family
\[
 \bigl\{\partial_t^n\partial_x^\alpha u:
 0\le n\le r,\ |\alpha|\le m\bigr\},
\]
using the displayed recurrence at the initial face. Higher-order persistence
realizes these fields as the actual derivatives of one velocity. Whenever
two finite constructions overlap in indices or time, local uniqueness
identifies their common coordinates. On the compact interval \([0,T]\),
persistence gives
\[
 \sum_{n=0}^r\|V_n\|_{L^2(0,T;H^m)}^2<\infty.
\]
Taking every finite \((r,m)\) proves
\eqref{f:eq:strict-independent-mixed-intersection}. The pressure formula and
Theorem~\ref{f:thm:formal-finite-row-pressure} place each selected pressure row
in its stated strict-horizon Sobolev space. Restriction compatibility makes
the resulting family one mixed jet of one velocity.
\end{proof}

This compact-horizon result identifies the rows of the one forced history.
The next theorem gives their stronger whole-terminal realization directly on
\(I=(0,T_*)\).  Restriction to \((0,T)\), followed by \(T\uparrow T_*\), is
a consequence of that realization.

For the whole-terminal step, keep this same forced history and fix finite
\((N,M)\).  The theorem below is the forced producing theorem.  Its analytic
input is the one prescribed triple \((u_0,\nu,f)\), its finite initial jet, and
the closed velocity--pressure--force graph.  It first supplies the adjacent
whole-terminal velocity memberships at that finite pair.  Their finite norm
sum then gives the bound uniform in every strict cutoff by restriction and
monotone convergence.  No unforced row, norm, history, or terminal value enters
this construction.

\begin{theorem}[Datum-generated compatible forced finite rectangles]
\label{f:thm:datum-finite-blocks}
Let \(\nu>0\), \(u_0\in\HsSigma(\R^3)\), and
\(f\in C^\infty([0,\infty);\mathcal S(\R^3;\R^3))\). Let
\((T_*,u,p)\) be the maximal pressure-complete classical history generated
by the single forced datum \((u_0,\nu,f)\). Assume, toward contradiction, that
\(T_*<\infty\), and put \(I=(0,T_*)\).  For finite \(N,M\), the corresponding
pressure-complete record is
\begin{equation}
 \cR_{N,M}[u_0,\nu,f;T_*]
 :=\left((V_n)_{n=0}^{N},(V_{n+1})_{n=0}^{N};
 (Q_n)_{n=0}^{N};(F_n)_{n=0}^{N}\right).
 \label{f:eq:forced-datum-rectangle-record}
\end{equation}
The forced construction produces, for \(0\le n\le N\),
\begin{equation}
 V_n\in L^2(I;H^{M+1}),\qquad
 V_{n+1}\in L^2(I;H^{M-1}).
 \label{f:eq:independent-adjacent-finite-rows}
\end{equation}
For \(0<S\le T_*\), define its velocity norm by
\begin{equation}
 \mathfrak R_{N,M}(u;S)
 :=\sum_{n=0}^{N}\left(
 \norm{V_n}_{L^2(0,S;H^{M+1})}^2
 +\norm{V_{n+1}}_{L^2(0,S;H^{M-1})}^2
 \right).
 \label{f:eq:forced-restricted-block-norm}
\end{equation}
The forced whole-terminal output is
\begin{equation}
 \sup_{0<S<T_*}\mathfrak R_{N,M}(u;S)
 =\mathfrak R_{N,M}(u;T_*)<\infty
 \qquad(N,M<\infty).
 \label{f:eq:forced-datum-block-output}
\end{equation}
Taking \(N\ge n\) and \(M=m\) gives the corresponding statement for every
finite pair \((n,m)\).
For \(N'\le N\) and \(M'\le M\), the larger block restricts to the same
velocity, pressure, and force coordinates as the smaller block:
\begin{equation}
 \left.\cR_{N,M}[u_0,\nu,f;T_*]\right|_{(N',M')}
 =\cR_{N',M'}[u_0,\nu,f;T_*].
 \label{f:eq:independent-intersection-record}
\end{equation}
Here restriction deletes the outer temporal rows and reads every retained
distribution in the lower Sobolev order.  The mixed coordinates are
\[
 J_{n,\alpha}=\partial_x^\alpha\partial_t^nu,
 \qquad
 \Pi_{n,\alpha}=\partial_x^\alpha\partial_t^np,
 \qquad
 F_{n,\alpha}=\partial_x^\alpha\partial_t^nf.
\]
For every finite \((n,m,\alpha)\),
\[
 J_{n,\alpha}\in L^2(I;H^m),\qquad
 \Pi_{n,\alpha}\in L^1(I;H^m),\qquad
 F_{n,\alpha}\in C([0,T_*];H^m).
\]
\end{theorem}

\begin{proof}[Datum construction on the complete forced graph]
Put \(I=(0,T_*)\), and fix finite \(N,M\).  The prescribed data first generate
the complete initial jet:
\[
 \begin{aligned}
 A_0&=u_0,\\
 P_n&=R_iR_j\sum_{a=0}^n\binom na(A_a)_i(A_{n-a})_j
      -(-\Delta)^{-1}\nabla\cdot F_n(0),\\
 A_{n+1}&=\nu\Delta A_n
 -\sum_{a=0}^n\binom na(A_a\cdot\nabla)A_{n-a}
 -\nabla P_n+F_n(0).
 \end{aligned}
\]
At step \(n\), this recurrence determines \(A_{n+1}\) from the preceding
velocity rows and the prescribed force row \(F_n(0)\).
Lemma~\ref{f:lem:formal-initial-jet} places every required initial row in every
finite spatial Sobolev order, and spatial differentiation supplies the initial
mixed coordinates.

For an equation indexed by \((n,\alpha)\), retain the adjacent temporal
coordinate \(J_{n+1,\alpha}\), the viscous coordinates
\(J_{n,\alpha+2e_\ell}\), and the pressure coordinates
\(\Pi_{n,\alpha+e_\ell}\), for \(1\le\ell\le3\), together with the direct
source coordinate \(F_{n,\alpha}\).  For every \(a\le n\),
\(\beta\le\alpha\), and velocity component \(i\), retain both transport
factors
\[
 J_{a,\beta,i},
 \qquad
 J_{n-a,\alpha-\beta+e_i}.
\]
Here \(e_i\) is the \(i\)-th spatial unit multi-index.  Set
\(\Gamma_\alpha:=\{\alpha\}\cup\{\alpha+e_\ell:1\le\ell\le3\}\).
For every \(\gamma\in\Gamma_\alpha\), every \(a\le n\), and every
\(\beta\le\gamma\), also retain the pressure-product factors
\(J_{a,\beta,i}\), \(J_{n-a,\gamma-\beta,j}\), and the pressure-source
coordinate \(F_{n,\gamma}\).  Each retained normalized pressure coordinate
then carries both its full velocity product and its prescribed source
potential:
\[
 \Pi_{n,\gamma}
 =R_iR_j\sum_{a=0}^n\sum_{\beta\le\gamma}
 \binom na\binom{\gamma}{\beta}
 J_{a,\beta,i}J_{n-a,\gamma-\beta,j}
 -(-\Delta)^{-1}\nabla\cdot F_{n,\gamma}
 \qquad(\gamma\in\Gamma_\alpha).
\]
The retained equation is therefore
\begin{equation}
 \mathscr E_{n,\alpha}(J,\Pi;F)
 :=J_{n+1,\alpha}
 +\sum_{a=0}^{n}\sum_{\beta\le\alpha}
  \binom na\binom{\alpha}{\beta}
  (J_{a,\beta}\cdot\nabla)J_{n-a,\alpha-\beta}
 +\nabla\Pi_{n,\alpha}
 -\nu\Delta J_{n,\alpha}-F_{n,\alpha}=0.
 \label{f:eq:forced-finite-equation-graph}
\end{equation}
There are finitely many retained coordinates for the fixed selected equations.
This dependency cover is not a second or truncated flow: every velocity and
pressure entry is a derivative of the one prescribed forced history, and
every force entry is a derivative of the one prescribed source.  Applying
\(\Pp\) reads the solenoidal evolution; applying \(I-\Pp\) reads the
normalized pressure equation.

The prescribed source adds one known affine coordinate to each coupled row.
At the block's velocity-row order, for every \(0\le n\le N\), Fourier
Cauchy--Schwarz and Young's inequality give the source payment
\begin{equation}
 2\left|\langle F_n,V_n\rangle_{H^M}\right|
 \le \nu\norm{V_n}_{H^{M+1}}^2
     +\nu^{-1}\norm{F_n}_{H^{M-1}}^2.
 \label{f:eq:forced-finite-row-source-payment}
\end{equation}
The first right-hand term is absorbed by the viscous
\(H^{M+1}\)-dissipation in the same row.  The second is finite on
\([0,T_*]\) because \(F_n\) is smooth in time with Schwartz spatial values.
For each retained mixed pressure index \(\gamma\in\Gamma_\alpha\), the
gradient component has the finite budget
\(\mathfrak F^{\rm p}_{n,\gamma,M}(I)\) from
\eqref{f:eq:force-pressure-budget}.  Thus the affine force and every
force-induced pressure coordinate are paid inside this finite forced record.

The initial face, complete dependency cover, two projections, and prescribed
affine source now form the input to the forced columnwise construction.
Applied to this one graph and this one history, it supplies, for every
\(n=0,\ldots,N\),
\[
 V_n\in L^2(I;H^{M+1}),
 \qquad
 V_{n+1}\in L^2(I;H^{M-1}).
\]
These are the memberships in
\eqref{f:eq:independent-adjacent-finite-rows}.  They are produced after the
force pairing and pressure-source payments above and before any terminal trace
or restart.

For this fixed finite pair, form the sum of their defining squared norms:
\[
 C^f_{N,M}[u_0,\nu,f;I]
 :=\sum_{n=0}^{N}\left(
  \norm{V_n}_{L^2(I;H^{M+1})}^2
  +\norm{V_{n+1}}_{L^2(I;H^{M-1})}^2
 \right)<\infty.
\]
Restriction to \((0,S)\) decreases each nonnegative integral.  Increasing
\(S\) to \(T_*\) and applying monotone convergence to the finite sum gives
\begin{equation}
 \sup_{0<S<T_*}\mathfrak R_{N,M}(u;S)
 =\mathfrak R_{N,M}(u;T_*)
 =C^f_{N,M}[u_0,\nu,f;I]<\infty.
 \label{f:eq:forced-uniform-cutoff-output}
\end{equation}
This proves \eqref{f:eq:forced-datum-block-output} in the construction's
actual order: prescribed triple, complete forced graph, force payments,
whole-terminal memberships, finite norm, restriction, monotone terminal
realization.  The scalar is the velocity norm of
\(\cR_{N,M}[u_0,\nu,f;T_*]\); the record also retains the pressure and
prescribed-force coordinates constrained by
\eqref{f:eq:forced-finite-equation-graph}.

Fix \(n,m,\alpha\).  Choosing \(N\ge n\) and
\(M\ge m+|\alpha|\) gives the stated
velocity and mixed-velocity memberships.  The normalized Riesz formula and
Theorem~\ref{f:thm:formal-finite-row-pressure} give the pressure membership,
including its prescribed source potential, while the prescribed force class
gives \(F_{n,\alpha}\in C([0,T_*];H^m)\).  The temporal recurrence, spatial
Leibniz rule, pressure normalization, and force jet commute with deletion of
the outer indices.  Overlap uniqueness identifies every retained coordinate
with the corresponding derivative of the original maximal history.  Hence
each enlarged rectangle restricts to the smaller one, so the finite records
satisfy the literal restriction identity
\eqref{f:eq:independent-intersection-record} on \(I\).
\end{proof}

The datum-generated blocks admit two all-orders readings.  Ranging over their
mixed indices gives the following intersection.  The next section independently
derives the completed-height identities and uses their coordinates to recover
the spatial column before the forced recurrence reconstructs the same mixed
rows.

\begin{corollary}[Mixed-index exhaustion of the forced finite rectangles]
\label{f:thm:datum-intersection}
Under the hypotheses of Theorem~\ref{f:thm:datum-finite-blocks},
\begin{equation}
 u\in\bigcap_{r,m\in\Nzero}H^r(I;H^m_\sigma(\R^3)),
 \qquad I=(0,T_*).
 \label{f:eq:independent-mixed-intersection}
\end{equation}
This is the direct mixed-index exhaustion of the datum-generated family
\(\bigl(\cR_{N,M}[u_0,\nu,f;T_*]\bigr)_{N,M\in\Nzero}\).
\end{corollary}

\begin{proof}

Fix finite \(r,m\).  A rectangle with temporal depth at least \(r\) and spatial
height at least \(m\) contains
\[
 \sum_{n=0}^{r}\norm{V_n}_{L^2(I;H^m)}^2<\infty.
\]
The recurrence identifies \(V_n\) with the weak derivative
\(\partial_t^nu\), which proves
\eqref{f:eq:independent-mixed-intersection}.  Choosing the spatial height higher
by \(|\alpha|\) gives every \(J_{n,\alpha}\).  The finite-row pressure estimate
gives every \(\Pi_{n,\alpha}\), and the prescribed force class gives every
\(F_{n,\alpha}\).  Equation~\eqref{f:eq:independent-intersection-record}
identifies these rows as literal common coordinates of every sufficiently
large block on this one forced history.
\end{proof}

\begin{definition}[Finite pressure-complete forced rectangle]
\label{f:def:finite-rectangle}
Let \(0<T<T_*\) and \(N,M\in\Nzero\).  Define
\[
 \cR_{N,M}[u_0,\nu,f;T]
 :=\left(
 (V_n|_{(0,T)})_{n=0}^{N},
 (V_{n+1}|_{(0,T)})_{n=0}^{N};
 (Q_n|_{(0,T)})_{n=0}^{N};
 (F_n|_{(0,T)})_{n=0}^{N}
 \right),
\]
with every row constrained by
Proposition~\ref{f:prop:formal-mixed-recurrences}.  Its associated scalar norms are
\begin{align}
 \mathfrak R_{N,M}(u;T)
 :={}&\sum_{n=0}^{N}
 \norm{V_n}_{L^2(0,T;H^{M+1})}^2
 +\sum_{n=0}^{N}
 \norm{V_{n+1}}_{L^2(0,T;H^{M-1})}^2,
 \label{f:eq:finite-rectangle-norm}\\
 \mathfrak Q_{N,M}(p;T)
 :={}&\sum_{n=0}^{N}
 \left(
 \norm{Q_n}_{L^1(0,T;H^M)}
 +\norm{\nabla Q_n}_{L^1(0,T;H^{M-1})}
 \right),
 \label{f:eq:finite-pressure-rectangle-norm}\\
 \mathfrak F_{N,M}(f;T)
 :={}&\sum_{n=0}^{N}\left(
 \norm{F_n}_{L^1(0,T;H^{M-1})}
 +\mathfrak F^{\rm p}_{n,0,M}(0,T)
 \right).
 \label{f:eq:finite-force-rectangle-norm}
\end{align}
\end{definition}

\begin{theorem}[Compatible forced rectangle norms]
\label{f:thm:datum-rectangles}
Under the hypotheses of Theorem~\ref{f:thm:datum-finite-blocks}, every finite
pair \((N,M)\) satisfies
\begin{equation}
 \mathfrak R_{N,M}^*[u_0,\nu,f,T_*]
 :=\sup_{0<T<T_*}\mathfrak R_{N,M}(u;T)<\infty.
 \label{f:eq:datum-rectangle-bound}
\end{equation}
For every finite \(n,m\), the adjacent coordinates obey
\begin{equation}
 V_n\in L^2(0,T_*;H^{m+1}),
 \qquad
 V_{n+1}\in L^2(0,T_*;H^{m-1}).
 \label{f:eq:datum-adjacent-rows}
\end{equation}
The constant in \eqref{f:eq:datum-rectangle-bound} may depend on the selected
finite orders.  Every entry is the corresponding derivative of the original
forced history.
\end{theorem}

\begin{proof}
Put \(I=(0,T_*)\).  By Theorem~\ref{f:thm:datum-finite-blocks}, for
\(0\le n\le N\),
\[
 V_n\in L^2(I;H^{M+1}),\qquad
 V_{n+1}\in L^2(I;H^{M-1}).
\]
Their finite squared sum
\[
 C^f_{N,M}:=\sum_{n=0}^{N}\left(
 \norm{V_n}_{L^2(I;H^{M+1})}^2
 +\norm{V_{n+1}}_{L^2(I;H^{M-1})}^2
 \right)
\]
is finite.  Restriction of each nonnegative integral gives
\[
 \mathfrak R_{N,M}(u;T)\le C^f_{N,M}
 \qquad(0<T<T_*),
\]
and proves \eqref{f:eq:datum-rectangle-bound}.  Selecting \((N,M)=(n,m)\)
gives \eqref{f:eq:datum-adjacent-rows}.  Enlarging the indices and deleting the
outer coordinates recovers the smaller block because every retained entry is
the same derivative of the prescribed forced history.
\end{proof}

The next section derives the critical-height and viscous relationships from
the same mixed jet.  It evaluates those identities on the same finite
blocks \(\cR_{N,M}[u_0,\nu,f;T_*]\), exhausts their completed heights through
the spatial Sobolev column, and returns through the forced recurrence to the
same mixed intersection.
